\documentclass[11pt]{amsart}

\usepackage[T1]{fontenc}
\usepackage{lmodern}
\usepackage{microtype}
\usepackage{amsmath,amssymb,amsthm}
\usepackage[hidelinks]{hyperref}

\hypersetup{
  pdftitle={Small-Parameter Counterexamples to Conjecture 3 of Larsson, Manabe, and Miyadera}
}

\newtheorem{theorem}{Theorem}
\newtheorem{lemma}[theorem]{Lemma}
\theoremstyle{remark}
\newtheorem{remark}[theorem]{Remark}

\newcommand{\mex}{\operatorname{mex}}
\newcommand{\G}{\mathcal{G}}

\title[Counterexamples to Conjecture 3]
{Small-Parameter Counterexamples to Conjecture 3 of Larsson, Manabe, and Miyadera}
\subjclass[2020]{91A46}
\keywords{Subtraction game, Nim with a pass, Sprague--Grundy function,
reverse mex}
\date{}

\begin{document}

\begin{abstract}
Larsson, Manabe, and Miyadera conjectured that three parameterized
families of subtraction Nim with a one-time pass satisfy several
properties of the Grundy sequence, among them two eventual mex
identities. Under the printed quantifiers $a,n\in\mathbb N$ in their
Conjecture~3, every positive integer $a$ yields a counterexample in
each subpart: the exceptional parameters we exhibit are respectively
$n=2$, $n=1$ and $n=2$. None of the three witnesses is covered by the
theorems they prove, which treat only the subtraction set
$\{2,4n,4n+2\}$ with $n\geq3$.
\end{abstract}

\maketitle

\section{Statement and notation}

Let $S\subset\mathbb N$ be a subtraction set. Write $\G_S(x,0)$ for
the Grundy value with no pass available and $\G_S(x,1)$ for the value
with a one-time pass available. Following the recursion used in
\cite{LMM},
\begin{align}
\G_S(x,0)
 &=\mex\{\G_S(x-s,0):s\in S,\ s\leq x\},
 \label{eq:no-pass}\\
\G_S(x,1)
 &=\mex\bigl(
     \{\G_S(x-s,1):s\in S,\ s\leq x\}
     \cup\{\G_S(x,0)\}
   \bigr),
 \label{eq:pass}
\end{align}
whenever $x>0$ admits a subtraction move, and $\G_S(0,1)=0$.
All positive positions used below admit the smallest subtraction, so
no convention concerning passes from positive terminal positions
enters the argument.

Conjecture~3 of \cite{LMM} concerns three parameterized subtraction
sets $S=\{s_1,s_2,s_3\}$ and the three properties displayed there as
Equations~(35), (36) and~(37), which we reproduce in that order:
\begin{align}
\G_S(x,0)
 &=\mex\{\G_S(x+s,0):s\in S\},
 \label{eq:reverse-no-pass}\\
\G_S(x,1)
 &=\mex\{\G_S(x-s,1):s\in S\},
 \label{eq:backward}\\
\G_S(x,1)
 &=\mex\{\G_S(x+s,1):s\in S\}.
 \label{eq:forward}
\end{align}
Its three subparts, labelled there (a), (b) and, in place of~(c),
``(iii)'', are each quantified over $a,n\in\mathbb N$, and each is a
conjunction:
\begin{enumerate}
\item[(a)] for $S=\{a,2an,2an+a\}$, that \eqref{eq:reverse-no-pass}
holds, that \eqref{eq:backward} and \eqref{eq:forward} hold for
$x\geq6an+(4a+1)$, and that the sequence of Grundy numbers enters into
a loop when $x\geq6an+(4a+1)$;

\item[(b)] for $S=\{a,(2n+1)a,(2n+3)a\}$, that
\eqref{eq:reverse-no-pass} holds, that \eqref{eq:backward} and
\eqref{eq:forward} hold for $x\geq a(2n+3)+1$, that $\G_S(x,1)<4$ for
every $x$, and that the sequence enters into a loop when
$x\geq a(2n+3)+1$;

\item[(iii)] for $S=\{a,(2n+1)a,(2n+5)a\}$, that
\eqref{eq:reverse-no-pass} holds, that $\G_S(x,1)<4$ for every $x$,
that \eqref{eq:backward} and \eqref{eq:forward} hold for
$x\geq a(2n+5)+1$, and that the sequence enters into a loop when
$x\geq a(2n+5)+1$.
\end{enumerate}
No restriction on $n$ beyond $n\in\mathbb N$ appears in the
conjecture.

Equation~\eqref{eq:forward} is the intended form of Equation~(37) of
\cite{LMM}, whose displayed left-hand side reads $\G(x)$ and so omits
the second argument. The sentence introducing that display reads
``(iii) The sequence of Grundy numbers $\G(x,1)$ has the reverse-mex
property'', which supplies the argument $1$. Reading the omitted
argument as $0$, and with it the arguments printed as $1$ on the right
of Equation~(37), would turn that equation into Equation~(35),
reproduced here as \eqref{eq:reverse-no-pass}, which each subpart
already asserts separately and with no threshold. The remaining, mixed
reading, with $\G(x,0)$ on the left and the printed argument $1$ on the
right, also fails at the witness of part~(a): at $a=1$ and $n=2$ the
table in the proof of Theorem~\ref{thm:main} gives $\mex\{2,1,0\}=3$
for the right-hand side at $x=18$, against a left-hand side of $0$.
The failure of
\eqref{eq:backward} recorded below does not depend on this reading, as
Equation~(36) is printed with both arguments throughout.

Since each subpart is a conjunction, it is refuted by refuting any one
conjunct; Theorem~\ref{thm:main} below refutes each subpart through the
conjunct \eqref{eq:backward}, and part~(a) also through
\eqref{eq:forward}. The other conjuncts are not contradicted at the
witnesses exhibited: \eqref{eq:reverse-no-pass} holds at each of them,
as the table in the proof of Theorem~\ref{thm:main} shows for part~(a)
and as the parity formula of Lemma~\ref{lem:odd} shows, through
Lemma~\ref{lem:scaling}, for the other two; and the bound
$\G_S(x,1)<4$ asserted in parts~(b) and~(iii) holds at every $x$ for
the odd sets covered by Lemma~\ref{lem:odd}, which give
$\G_S(x,1)\leq2$. Nothing is asserted here about the eventual
periodicity claimed alongside.

\begin{lemma}[Scaling]\label{lem:scaling}
Suppose $1\in S$ and $a\in\mathbb N$. Then, for every
$k\in\mathbb Z_{\geq0}$ and $p\in\{0,1\}$,
\[
\G_{aS}(ak,p)=\G_S(k,p),
\qquad
aS:=\{as:s\in S\}.
\]
\end{lemma}

\begin{proof}
Induct on $k$, proving the case $p=0$ before the case $p=1$ at each
$k$. For $p=0$, every option from $ak$ is of the form $a(k-s)$, so
\eqref{eq:no-pass} and the induction hypothesis give the claim.
Having established the no-pass value at $ak$, Equation~\eqref{eq:pass}
gives the with-pass value from the lower with-pass values and the
current no-pass value. These option sets coincide with those at $k$
for $S$. The initial values at $k=0$ are both zero, and $1\in S$
ensures that every positive position $ak$ admits the subtraction $a$.
\end{proof}

\begin{lemma}[Odd subtraction sets]\label{lem:odd}
Let $S$ consist entirely of odd integers and contain $1$. Then
\[
\G_S(x,0)=
\begin{cases}
0,&x\text{ even},\\
1,&x\text{ odd},
\end{cases}
\]
and
\[
\G_S(x,1)=
\begin{cases}
0,&x=0,\\
1,&x>0\text{ even},\\
2,&x\text{ odd and }x\in S,\\
0,&x\text{ odd and }x\notin S.
\end{cases}
\]
\end{lemma}

\begin{proof}
Both formulas are established by a single simultaneous strong
induction on $x$, the parity formula for $\G_S(x,0)$ being proved
first at each $x$, since the with-pass value at $x$ uses the no-pass
value at $x$. The first formula follows because every legal move
changes parity. For $x>0$ even, all subtraction options have odd heap
size, so their with-pass values are $0$ or $2$ by induction, while
the pass contributes $\G_S(x,0)=0$; hence the mex is $1$.

For $x$ odd, every subtraction option has even heap size. If $x\in S$,
one option is $0$, while every other legal subtraction option and the
pass have value $1$, so the mex is $2$. If $x\notin S$, all subtraction
options and the pass have value $1$, so the mex is $0$.
\end{proof}

\section{Counterexamples}

\begin{theorem}\label{thm:main}
Under the printed quantifiers $a,n\in\mathbb N$, every subpart of
Conjecture~3 in \cite{LMM} is false. More precisely, for every
$a\in\mathbb N$:
\begin{enumerate}
\item
In part~(a), take $n=2$, so $S=a\{1,4,5\}$. At $x=18a$, both
\eqref{eq:backward} and \eqref{eq:forward} fail.

\item
In part~(b), take $n=1$, so $S=a\{1,3,5\}$. At $x=6a$,
Equation~\eqref{eq:backward} fails.

\item
In the third subpart, printed as ``(iii)'', take $n=2$, so
$S=a\{1,5,9\}$. At $x=10a$, Equation~\eqref{eq:backward} fails.
\end{enumerate}
\end{theorem}

\begin{proof}
By Lemma~\ref{lem:scaling}, it is enough to compute at scale $a=1$.

For $T=\{1,4,5\}$, successive use of \eqref{eq:no-pass} and
\eqref{eq:pass} gives
\[
\begin{array}{c|rrrrrrrrrrrr}
x&0&1&2&3&4&5&6&7&8&9&10&11\\ \hline
\G_T(x,0)&0&1&0&1&2&3&2&3&0&1&0&1\\
\G_T(x,1)&0&2&1&0&1&4&0&2&3&0&1&3
\end{array}
\]
and
\[
\begin{array}{c|rrrrrrrrrrrr}
x&12&13&14&15&16&17&18&19&20&21&22&23\\ \hline
\G_T(x,0)&2&3&2&3&0&1&0&1&2&3&2&3\\
\G_T(x,1)&0&1&3&0&1&2&4&2&3&0&1&0.
\end{array}
\]
Consequently,
\[
\mex\{\G_T(17,1),\G_T(14,1),\G_T(13,1)\}
 =\mex\{2,3,1\}
 =0
 \neq4
 =\G_T(18,1),
\]
and
\[
\mex\{\G_T(19,1),\G_T(22,1),\G_T(23,1)\}
 =\mex\{2,1,0\}
 =3
 \neq4
 =\G_T(18,1).
\]
Thus both required identities fail in part~(a).

For $S=\{1,3,5\}$, Lemma~\ref{lem:odd} gives
\[
\G_S(1,1)=\G_S(3,1)=\G_S(5,1)=2,
\qquad
\G_S(6,1)=1.
\]
Hence
\[
\mex\{\G_S(5,1),\G_S(3,1),\G_S(1,1)\}
 =\mex\{2,2,2\}
 =0
 \neq1
 =\G_S(6,1).
\]
This disproves part~(b).

For $U=\{1,5,9\}$, Lemma~\ref{lem:odd} similarly gives
\[
\G_U(1,1)=\G_U(5,1)=\G_U(9,1)=2,
\qquad
\G_U(10,1)=1,
\]
so
\[
\mex\{\G_U(9,1),\G_U(5,1),\G_U(1,1)\}
 =\mex\{2,2,2\}
 =0
 \neq1
 =\G_U(10,1).
\]
This disproves the third subpart.

Finally, the conjectured thresholds specialize at the chosen
parameters: $6an+(4a+1)=16a+1$ at $n=2$ in part~(a),
$a(2n+3)+1=5a+1$ at $n=1$ in part~(b), and $a(2n+5)+1=9a+1$ at $n=2$
in the third subpart. The witnesses satisfy
\[
18a\geq16a+1,
\qquad
6a\geq5a+1,
\qquad
10a\geq9a+1,
\]
for every $a\geq1$. In part~(b) and in the third subpart the witness is
$x=\min S+\max S$, which exceeds the corresponding threshold by exactly
$a-1$; the two witnesses therefore meet their thresholds with equality
at $a=1$.
\end{proof}

\begin{remark}[Scope]
The abstract and the principal definitions of \cite{LMM} restrict the
family $\{2,4n,4n+2\}$ to $n\geq3$, the abstract adding: ``We do not
treat the case that $n=1$ or $n=2$ in this article.'' Conjecture~3,
however, separately quantifies $a,n\in\mathbb N$ in each subpart, and
Theorem~\ref{thm:main} addresses that printed statement. None of the
three witnesses lies in the family treated by the theorems of
\cite{LMM}. In part~(a) the set $\{a,2an,2an+a\}$ has the form
$\{2,4m,4m+2\}$ only when $a=2$, and there the choice $n=2$ gives
$\{2,8,10\}$, one of the two cases explicitly set aside. In part~(b)
and in the third subpart no member of the family can have that form:
matching the least elements forces $a=2$, and the middle element is
then $4n+2\equiv2\pmod4$, whereas $4m\equiv0\pmod4$. Accordingly
Theorem~\ref{thm:main} does not contradict the proved $n\geq3$ results
in \cite{LMM}, and it
makes no assertion about Conjecture~3 after imposing the additional
restriction $n\geq3$. The parameters of Theorem~\ref{thm:main} are not
claimed to be the only exceptional ones. For $a=1$, where $1\in S$ and
no positive position is terminal, direct computation for
$n=1,\dots,12$ finds no others: \eqref{eq:backward} or
\eqref{eq:forward} fails above the stated threshold at $n=2$, $n=1$
and $n=2$ respectively, and at no other $n$ in that range. For
$a\geq2$ each family contains positions from which no subtraction is
legal, and the set of exceptional $n$ then depends on the convention
adopted for the pass at such positions. Theorem~\ref{thm:main} is
independent of that convention, since every position it uses admits
the subtraction $a$.
\end{remark}

\begin{thebibliography}{1}

\bibitem{LMM}
U.~Larsson, H.~Manabe, and R.~Miyadera,
\emph{A subtraction Nim with a pass},
arXiv:2605.14321v1 [math.CO] (2026),
\href{https://arxiv.org/abs/2605.14321}{arXiv:2605.14321}.

\end{thebibliography}

\end{document}